\begin{table}[H]
\centering
\caption{Variables del VAR y sus transformaciones}
\label{tab:transformaciones}
\small
\begin{tabular}{@{}llll@{}}
\toprule
Variable & Tipo & Transformación & Unidad \\
\midrule
ISE & Endógena 1 & $100 \times (\log ISE_t - \log ISE_{t-1})$ & \% mensual \\
Inflación & Endógena 2 & $100 \times (\log IPC_t - \log IPC_{t-1})$ & \% mensual \\
Expectativas & Endógena 3 & $100 \times$ mediana de la expectativa a 12 meses & \% \\
TPM & Endógena 4 & Promedio mensual de la serie diaria & \% \\
ITCR & Endógena 5 & $100 \times (\log ITCR_t - \log ITCR_{t-1})$ & \% mensual \\
Brent & Exógena & $100 \times (\log Brent_t - \log Brent_{t-1})$, en $t$ y $t-1$ & \% mensual \\
Tasa Fed & Exógena & Nivel, en $t$ y $t-1$ & \% \\
Estacionales & Exógena & 11 dummies, una por mes de febrero a diciembre & --- \\
Puntuales & Exógena & Una dummy por mes: mar--dic 2020 y may 2021 & --- \\
\bottomrule
\end{tabular}
\notascuadro{\textit{Nota.} El orden de las endógenas es el del VAR. Un aumento del ITCR es una depreciación real del peso.

Fuente: Elaboración propia con datos de DANE, Banco de la República y FRED.}
\end{table}
